\section*{Context}

The global energy landscape is undergoing significant changes, driven by the urgent need to reduce greenhouse gas emissions and transition toward sustainable energy sources. Solar photovoltaic (PV) technology has become one of the central pillars of this transition. By the end of 2025, cumulative installed PV capacity reached approximately 2,383~GW worldwide \citep{irena2026}. This rapid growth has been enabled by continuous reductions in module costs, improvements in conversion efficiency, and international climate commitments following the Paris Agreement.

PV systems offer several characteristics that make them particularly attractive for large-scale deployment. Their modular architecture allows installation at scales ranging from residential rooftops to utility-scale power plants. They operate silently, produce no direct greenhouse gas emissions during electricity generation, and contain no rotating machinery, resulting in relatively low maintenance requirements. Utility-scale PV installations typically incur operation and maintenance costs of around 13~USD/kW/year, compared with 20--93~USD/kW/year for onshore wind systems \citep{irena_costs2025}. Furthermore, service lifetimes exceeding 25 years are routinely achieved in well-maintained installations.

Algeria possesses exceptional solar resources and considerable potential for photovoltaic development. In recent years, the country has accelerated investments in solar energy as part of its strategy to diversify the national energy mix and reduce dependence on fossil fuels. Sonelgaz currently operates 29 PV plants with a combined installed capacity of 383.1~MW, producing 590.44~GWh in 2024. In addition, 21 new plants representing approximately 3,200~MWc of capacity are currently under construction \citep{energygov2025}. Maximizing the return on these investments requires not only successful deployment but also reliable operation throughout the systems' service lifetimes. Effective fault detection and maintenance strategies are therefore becoming increasingly important.

\section*{Problem Statement and Motivation}

The economic and operational benefits of photovoltaic systems depend on their ability to maintain expected performance over long periods of operation. In practice, however, PV installations are exposed to various fault conditions that progressively degrade system performance. Unlike conventional power generation technologies, where failures are often mechanically audible or visually apparent, PV faults are generally invisible to direct observation and can only be inferred from deviations in electrical and environmental measurements.

When faults remain undetected, energy production decreases, equipment degradation accelerates, maintenance costs increase, and system reliability deteriorates. In severe cases, electrical anomalies may also create safety hazards. These consequences are particularly significant for utility-scale PV plants, where even small performance losses can translate into substantial economic impacts over time.

Traditional maintenance approaches rely on periodic inspections performed by qualified technicians. While such methods may be suitable for small installations, they become increasingly impractical as plant size grows. Modern utility-scale facilities may contain tens of thousands of modules distributed across hundreds of hectares, making frequent manual inspection both costly and operationally inefficient. Inspection intervals measured in weeks or months are often incompatible with the timescale at which faults develop and propagate.

Automated Fault Detection and Diagnosis (FDD) systems address this challenge by continuously monitoring operational data and identifying abnormal behaviour in real time. Recent advances in machine learning and deep learning have demonstrated promising results for fault detection and classification in PV systems. However, two important limitations remain. First, a large proportion of existing studies rely on simulated datasets that do not fully represent the complexity and variability of real-world operating conditions. Second, relatively few works consider the computational constraints associated with deployment on embedded edge devices. Addressing these challenges requires not only accurate diagnostic models but also a rigorous end-to-end methodology encompassing data preparation, feature engineering, model validation, and practical deployment.

\section*{Objectives}

This thesis aims to develop a reproducible, lightweight, and operationally relevant FDD framework for photovoltaic systems. The work is structured around three main objectives:

\begin{enumerate}
    \item \textbf{Rigorous data pipeline.} Design and implement a complete workflow covering data ingestion, temporal splitting, preprocessing, and feature engineering. Particular attention is given to handling imperfections in real-world measurements and preventing data leakage throughout model development.

    \item \textbf{Lightweight FDD models.} Propose and evaluate efficient machine learning and deep learning architectures for two complementary diagnostic tasks: anomaly detection, which determines whether a fault condition exists, and fault classification, which identifies the corresponding fault type. Beyond diagnostic performance, the models are assessed according to their suitability for deployment on resource-constrained hardware.

    \item \textbf{Edge deployment.} Deploy the validated framework on a Jetson Nano embedded platform and develop a graphical user interface for real-time monitoring, thereby bridging the gap between laboratory evaluation and practical field implementation.
\end{enumerate}